\chapter{Édition des rapports}
\minitoc

\section{Le principe}

Rien n'est produit à l'ouverture de l'onglet \textbf{Rapports}. C'est
délibéré : un rapport que personne n'a demandé n'a pas d'objet, et en produire
un d'office ferait croire que la plateforme ne sait en faire qu'un seul.

Vous choisissez donc deux choses : \textbf{ce que le rapport doit couvrir}, et
\textbf{sous quel format} l'obtenir.

\section{Choisir le périmètre}

\begin{table}[h]
\centering
\begin{tabular}{@{}p{4.2cm}p{9.6cm}@{}}
\toprule
\textbf{Périmètre} & \textbf{Ce qu'il retient} \\
\midrule
Une période &
Toute l'activité sur 24 heures, 7, 30, 90 jours ou un an. \\[6pt]
Une intervention &
Le compte rendu détaillé d'une seule affaire, à annexer à un courrier ou à une
procédure. \\[6pt]
Une famille d'attaque &
Réseau, applicative, comportementale ou infrastructure. \\[6pt]
Un niveau de gravité &
Le niveau retenu \textbf{et au-dessus} : demander les incidents graves en
espérant que les plus graves en soient absents n'aurait pas de sens. \\[6pt]
Un type du catalogue &
Un seul des types métier, appris compris. \\
\bottomrule
\end{tabular}
\caption{Les cinq périmètres}
\end{table}

Un bouton \textbf{Prévisualiser} affiche le document à l'écran, dans la mise en
page exacte qu'aura le document imprimé. On télécharge ensuite, si le résultat
convient.

\section{Choisir le format}

\begin{table}[h]
\centering
\begin{tabular}{@{}p{2.4cm}p{4cm}p{7.4cm}@{}}
\toprule
\textbf{Format} & \textbf{Produit par} & \textbf{Usage} \\
\midrule
PDF & ReportLab &
Document officiel, prêt à signer et à classer. \\[6pt]
Word & python-docx &
Reprise et annotation : styles de titre et tableaux réels, pas une image. \\[6pt]
Markdown & --- &
Texte lisible sans aucun outil, à coller dans un courriel. \\[6pt]
JSON & bibliothèque standard &
Transmission à un autre système. Les tableaux restent des tableaux. \\
\bottomrule
\end{tabular}
\caption{Les quatre formats}
\end{table}

\begin{tcolorbox}[colback=green!4!white, colframe=darkgreen,
  title=\textbf{Bon à savoir --- les quatre versions disent la même chose},
  fonttitle=\bfseries, boxrule=0.6pt]
Le document est composé \textbf{une fois}, sous forme de blocs, puis rendu dans
les quatre formats. Écrire quatre générateurs produirait immanquablement des
écarts, et un rapport officiel qui ne dit pas la même chose selon le fichier
ouvert n'est plus opposable. Un test compare les quatre à chaque
modification.
\end{tcolorbox}

\section{Ce que contient un rapport}

Le document reprend la présentation des actes administratifs camerounais :
titulature bilingue en trois colonnes — français à gauche, emblème de l'Agence
au centre, anglais à droite — référence d'ordre, objet, parties numérotées en
chiffres romains, lieu, date et bloc de signature.

Un rapport périodique comporte neuf parties : objet, conditions de
fonctionnement, vue d'ensemble chiffrée, nature des menaces traitées avec
diagrammes, interventions conduites, gestes appliqués, situations où la
plateforme s'est abstenue, traçabilité, appréciation d'ensemble.

Un compte rendu d'intervention en comporte six : identification, ce qui a été
observé, ce que la plateforme a décidé, gestes appliqués, déroulement
chronologique, conclusion.

\section{La numérotation des documents}

Chaque rapport porte une référence de la forme :

\begin{center}
\texttt{N° 4503/RAP/MINPOSTEL/ANTIC/CIRT/2026}
\end{center}

Elle est calculée à partir du périmètre et du jour d'édition. Deux éditions du
même rapport le même jour portent \textbf{le même numéro} — ce qu'attend un
service d'archives : rééditer un rapport ne crée pas une pièce nouvelle.

\section{Le langage employé}

Aucun identifiant technique n'atteint la page. Une table de correspondance
unique traduit tout le vocabulaire :

\begin{center}
\begin{tabular}{@{}ll@{}}
\toprule
\textbf{Dans le code} & \textbf{Dans le rapport} \\
\midrule
\texttt{firewall:block\_ip} & Blocage d'une adresse au pare-feu \\
\texttt{edr:isolate\_host} & Mise en quarantaine réseau d'une machine \\
\texttt{no\_grounded\_context} & Aucune réponse engagée --- la menace n'est
pas documentée \\
\texttt{rolled\_back} & Annulée par la plateforme \\
\bottomrule
\end{tabular}
\end{center}

Le motif d'une décision est \textbf{recomposé} à partir des données
structurées, jamais recopié du moteur : celui-ci cite des identifiants de
playbook et une empreinte de politique, qui ont leur place dans le registre,
pas dans un document que lit un directeur.

\begin{tcolorbox}[colback=green!4!white, colframe=darkgreen,
  title=\textbf{Bon à savoir}, fonttitle=\bfseries, boxrule=0.6pt]
Un test refuse d'avance tout geste ajouté au catalogue sans traduction. Sans
lui, un geste nouveau finirait par apparaître sous son identifiant technique
dans un document officiel, et personne ne s'en apercevrait avant l'impression.
\end{tcolorbox}
